\chapter*{本册结构}
\addcontentsline{toc}{chapter}{本册结构}

第一册围绕一条基础链路展开：

\begin{center}
\begin{minipage}{0.86\linewidth}
\begin{verbatim}
C/C++ 源码
  -> 编译、汇编、链接
  -> 可执行文件和进程
  -> x86-64 指令、寄存器、内存和控制流
  -> System V AMD64 ABI
  -> C++ 与汇编互操作
  -> 可信性能测量
\end{verbatim}
\end{minipage}
\end{center}

这条链路决定了本册的三部分结构。

\section*{底层模型}

第一部分从程序如何成为进程讲起，再进入 CPU 执行、数据表示、对象布局、指针、虚拟地址空间、Linux 工具和调试证据。这里的重点不是记住工具命令，而是建立一组可验证事实：运行的是哪份二进制，指令改变了哪些架构状态，对象在内存中如何排列，工具输出能证明什么。

\section*{汇编与 ABI}

第二部分把前面的机器状态具体化为 x86-64 汇编。整数指令、寻址模式、flags、条件跳转、循环、函数调用和 \code{switch} 都要同时从 C++ 语义、汇编形态和工具证据三侧解释。System V AMD64 ABI 负责把函数边界变成二进制合同：参数从哪里进入，返回值从哪里出去，栈如何对齐，哪些寄存器由谁保存。C++ 与汇编互操作则把这些规则放进可维护工程：public wrapper、internal C ABI、汇编 kernel、测试、反汇编审计和 benchmark 必须分层。

\section*{可信测量}

第三部分处理性能数字。运行时间本身只是观察，不是结论。可信 benchmark 必须先证明 correctness，再固定输入、计时区域、重复次数、统计口径、环境和二进制身份。第 15 章建立一次测量的基本协议，第 16 章把协议扩展为长期 benchmark suite，使结果可以复现、比较和维护。

\section*{进入第二册}

第二册会继续展开 cache/TLB、现代乱序核心、SIMD、HPC kernel、线程、NUMA、PMU 和 AI CPU 算子。第一册不提前浅讲这些主题，而是先把它们共同依赖的基础打通：源码语义、对象布局、二进制身份、汇编阅读、ABI 边界和测量证据。没有这些基础，高级优化很容易退化成口号；有了这些基础，后续主题才能落到可以验证的机器事实上。
